\chapter*{术语表}
\phantomsection
\addcontentsline{toc}{chapter}{术语表}

\begin{center}
\small
\begin{tabular}{@{}p{0.24\linewidth}p{0.66\linewidth}@{}}
\toprule
术语 & 实践含义 \\
\midrule
manifest & 模型资产清单，记录 tensor name、dtype、shape、offset、checksum 和 tokenizer/配置版本。 \\
shape verifier & 在 kernel 运行前检查 shape、dtype、layout、量化 descriptor 和 state edge 是否匹配。 \\
reference decoder & 以清楚正确为目标的 decoder 实现，用来生成 trace 和 golden。 \\
oracle & 比较候选实现与 reference 的判定器，可使用精确相等或误差界。 \\
trace checkpoint & 推理链路中的可复查节点，至少包含 shape、stride、dtype、layout、checksum 和摘要值。 \\
KV cache & attention 历史 K/V 状态，跨 token 存活，不能被普通 activation planner 复用。 \\
packing & 把权重改写成更适合热循环和 SIMD 的布局。 \\
backend & 执行计划落到机器的实现，例如 scalar、AVX2、AVX-512、GPU 或库调用。 \\
fallback & 后端能力不满足时显式降级到可验证路径，并记录原因。 \\
SLO & 服务目标，例如 TTFT、TPOT、queue wait、拒绝率和取消回收时间。 \\
\bottomrule
\end{tabular}
\end{center}
